\documentclass[10pt,a4paper]{amsart}
\usepackage[margin=19mm]{geometry}
\usepackage{amsmath,amssymb,mathtools}
\usepackage[scr=rsfs,cal=euler]{mathalfa}
\usepackage{xfrac}
\usepackage[unicode,hidelinks,bookmarks=false]{hyperref}
\newcommand{\ie}{i.e.\ }
\newcommand{\cf}{cf.\ }
\newcommand{\resp}{respectively\ }
\newcommand{\Subsection}{\subsection}
\providecommand{\nobreakdash}{\nobreak\hbox{-}}
\setlength{\emergencystretch}{2em}
\multlinegap0pt
\usepackage{amssymb}
\usepackage{mathtools}
\usepackage[shortlabels]{enumitem}
\usepackage{xcolor}
\usepackage{tikz}
\newtheorem{lemma}{Lemma}
\DeclareMathOperator{\Ad}{Ad}
\DeclareMathOperator{\Mat}{M}
\DeclareMathOperator{\diag}{diag}
\DeclareMathOperator{\rk}{rk}
\title{Uniqueness of Rank-Metric Completions of Bratteli Systems}
\author{Baojie Jiang}
\date{Source: arXiv:2609.12924v1 (CC BY 4.0)}
\begin{document}
\maketitle

\noindent\emph{Source and licence.} Uniqueness of Rank-Metric Completions of Bratteli Systems. Version arXiv:2609.12924v1 (CC BY 4.0), distributed by arXiv under the Creative Commons Attribution 4.0 International licence, http://creativecommons.org/licenses/by/4.0/, as recorded in the arXiv licence metadata for this exact version and shown on the abstract page https://arxiv.org/abs/2609.12924v1. Full licence notice: \texttt{licenses/arxiv-2609.12924-CC-BY-4.0.txt}.\par\smallskip

\noindent\textit{Editorial scope.} Counted unit: the article's lemma lem:finite-model-estimates (source0.tex lines 2265--2298). The article's complete original proof of it is delivered with it (source0.tex lines 2300--2475, one \texttt{proof} environment of 176 lines). No mathematics is written, repaired or re-derived: the statement and the proof are the article's own lines, in the article's own notation, and no retained line is substituted or re-flowed.  Retained uncounted as the source-local context the proof needs: lines 2124--2133; lines 2141--2144; lines 2156--2159; lines 2223--2227.  Nothing was omitted from any retained span, so the package carries no omission record.  No retained line contains a citation, so the wrapper carries no bibliography and no bibliography entry.  No retained line contains a figure environment, an image-inclusion command or drawing code, so no image is delivered and the package declares no asset dependency.  Editorial formatting only, in addition: an AMS \texttt{amsart} preamble that restates the article's own declarations and macros used by the retained lines, namely the environments \texttt{lemma} on separate counters, together with the standard packages those lines need.  The retained environments print in this document as Lemma 1 in order of appearance, whereas the article prints the same items with the numbering of its own full document.  The complete native source of the article as distributed in the arXiv e-print for this exact version is bundled unchanged as \texttt{sources/210162/source0.tex}.
\begin{equation}\label{eq:alignment-coordinate-copies}
 \begin{aligned}
 (\phi_{n,m}(x))_w
 &=\sum_{v\in V_m}\sum_{a=1}^{a_{w,v}^{(n,m)}}
       J_{w,v,a}x_vJ_{w,v,a}^{T},\\
 \sigma_{n,q,w}^0\rho_{m,q}(x)
 &=\sum_{v\in V_m}\sum_{b=1}^{t_wc_{m,q}(v)}
       K_{w,v,b}x_vK_{w,v,b}^{T}.
 \end{aligned}
\end{equation}

\begin{equation}\label{eq:alignment-coordinate-match}
 U_wK_{w,v,a}=J_{w,v,a},
 \qquad \forall v\in V_m,\ \forall 1\leq a\leq \min\{a_{w,v}^{(n,m)},t_wc_{m,q}(v)\}.
\end{equation}

\begin{equation}\label{eq:aligned-packing-P}
 \sigma_{m,n,q}^{U,P}
 :=\sigma_{m,n,q}^{U}\Ad(P^{-1}).
\end{equation}

\begin{equation}\label{eq:return-map-composition}
 \rho_{m,n,q,r}^{U,P}\sigma_{m,n,q}^{U,P}(y)
 =\diag(1_{D_{n,q,r}}\otimes y,0_{(r-D_{n,q,r})q}),
 \qquad \forall y\in\Mat_q(R).
\end{equation}

\begin{lemma}
\label{lem:finite-model-estimates}
Fix $1\leq m<n$, $q,r\geq1$, a permutation matrix $P\in\Mat_q(R)$, and an
$(m,n,q)$-alignment $U$.
Then
\begin{equation}\label{eq:eta}
 0\leq\eta_{m,q}\leq\frac1q\sum_{v\in V_m}p_m(v).
\end{equation}
For every $d\geq1$,
\begin{equation}\label{eq:amplified-eta}
 \operatorname{dist}_d(\rho_{m,q}^{P})=d\eta_{m,q},
 \qquad
 \operatorname{dist}_d(\sigma_{m,n,q}^{U,P})=d\tau_{n,q},
 \qquad
 \operatorname{dist}_d(\rho_{m,n,q,r}^{U,P})=d\eta_{n,qr}.
\end{equation}
For every $X\in\Mat_d(A_m)$,
\begin{equation}\label{eq:parameter-first-square}
 \rk_{\alpha,n}\bigl(
   \phi_{n,m}(X)-
   \sigma_{m,n,q}^{U,P}\rho_{m,q}^{P}(X)
 \bigr)
 \leq d\bigl(\xi_{m,n}+\eta_{m,q}+\tau_{n,q}\bigr).
\end{equation}
For every $Y\in\Mat_d(\Mat_q(R))$,
\begin{equation}
   \label{eq:parameter-second-square}
  \rk_{qr}\bigl(
   \jmath_{q,r}(Y)-
   \rho_{m,n,q,r}^{U,P}\sigma_{m,n,q}^{U,P}(Y)
 \bigr)
 \leq d\bigl(\eta_{n,qr}+\tau_{n,q}\bigr).
\end{equation}
\end{lemma}

\begin{proof}
The floor inequality gives, for every $v\in V_m$,
  \[
\frac{c_{m,q}(v)p_m(v)}q\leq\alpha_m(v)<\frac{p_m(v)}q + \frac{c_{m,q}(v)p_m(v)}q,
\]
Equivalently,
\[
0\leq\alpha_m(v) -  \lambda_{m,q}(v)<\frac{p_m(v)}q.
\]
Moreover,
\[
  \eta_{m,q}
  = 1- \sum_{v\in V_m}\lambda_{m,q}(v)
 =\sum_{v\in V_m}\bigl(\alpha_m(v)-\lambda_{m,q}(v)\bigr)\leq\frac1q\sum_{v\in V_m}p_m(v).
\]

For $X\in\Mat_d(A_m)$, by definition we have
\[
  \rho_{m,q}^{P}(X)
  =(1_d\otimes P)\rho_{m,q}(X)(1_d\otimes P^{-1}).
\]
Block additivity and permutation invariance give
\[
 \begin{aligned}
 \rk_q(\rho_{m,q}^{P}(X))
   & =  \rk_q (\rho_{m,q}(X))\\
   & =\frac1q\sum_{v\in V_m}c_{m,q}(v)\rk(X_v)\\
 &=\sum_{v\in V_m}\frac{c_{m,q}(v)p_m(v)}q
        \rk_{p_m(v)}(X_v)\\
 &=\sum_{v\in V_m}\lambda_{m,q}(v)\rk_{p_m(v)}(X_v).
 \end{aligned}
\]
Since
\[
\rk_{\alpha,m}(X)-\rk_q(\rho_{m,q}^{P}(X))=\sum_v(\alpha_m(v)-\lambda_{m,q}(v))\rk_{p_m(v)}(X_v),
\]
we have
\[
0\leq \rk_{\alpha,m}(X)-\rk_q(\rho_{m,q}^{P}(X))\leq d\eta_{m,q}.
\]
The first inequality follows from
$0\leq\lambda_{m,q}(v)\leq\alpha_m(v)$, and the second from
$0\leq\rk_{p_m(v)}(X_v)\leq d$.
At $X=1_d$, every component satisfies $\rk_{p_m(v)}(X_v)=d$,
so the bound is attained.  Thus
\[
 \operatorname{dist}_d(\rho_{m,q}^{P})=d\eta_{m,q}.
\]
Since $\rho_{m,n,q,r}^{U,P}=\rho_{n,qr}^{P'}$, the same calculation
with $(n,qr,P')$ gives, for $Z\in\Mat_d(A_n)$,
\[
  \sum_{w\in V_n}(\alpha_n(w)-\lambda_{n,qr}(w))
  \rk_{p_n(w)}(Z_w)\leq d\eta_{n,qr}.
\]
Again $Z=1_d$ attains equality, so
\[
\operatorname{dist}_d(\rho_{m,n,q,r}^{U,P})=d\eta_{n,qr}.
\]

For $Y\in\Mat_d(\Mat_q(R))$, precomposition by $\Ad(P^{-1})$
preserves its rank.
After flattening and conjugation by a permutation matrix independent of $Y$,
the $w$-component of $\sigma_{m,n,q}^{U,P}(Y)$ is block diagonal with
$t_w$ copies of $Y$ and a zero block.  Hence
\[
 \rk_{\alpha,n}(\sigma_{m,n,q}^{U,P}(Y))
 =\sum_w\frac{\alpha_n(w)}{p_n(w)}t_w\rk(Y)
 =\rk_q(Y)\sum_w\alpha_n(w)\frac{t_wq}{p_n(w)}.
\]
Thus,
\[
\rk_q(Y) - \rk_{\alpha,n}(\sigma_{m,n,q}^{U,P}(Y)) =
\rk_q(Y)\left(1-\sum_w\alpha_n(w)\frac{t_wq}{p_n(w)}\right) =
\rk_q(Y)\sum_w\alpha_n(w) (1-\theta_w),
\]
the last equality using $\sum_w\alpha_n(w)=1$.
Since $0\leq\rk_q(Y)\leq d$, this difference lies in $[0,d\tau_{n,q}]$, and the bound is
attained at $Y=1_d\in\Mat_d(\Mat_q(R))$.
This proves
\[
\operatorname{dist}_d(\sigma_{m,n,q}^{U,P})=d\tau_{n,q}.
\]

For \eqref{eq:parameter-first-square}, cancellation of
$\Ad(P^{-1})\Ad(P)$ in \eqref{eq:aligned-packing-P} reduces the composition to
$\sigma_{m,n,q}^{U}\rho_{m,q}$.
Put $\widehat J_{v,a}=1_d\otimes J_{w,v,a}$ and
$\widehat K_{v,b}=1_d\otimes(U_wK_{w,v,b})$ for each fixed $w$.
By \eqref{eq:alignment-coordinate-match},
$\widehat J_{v,a}=\widehat K_{v,a}$ for
$1\leq a\leq\min\{a_{w,v}^{(n,m)},t_wc_{m,q}(v)\}$.
Hence \eqref{eq:alignment-coordinate-copies}, applied entrywise, yields
\[
 \begin{aligned}
 &(\phi_{n,m}(X)-\sigma_{m,n,q}^{U}\rho_{m,q}(X))_w\\
 &=\sum_v\sum_{t_wc_{m,q}(v)<a\leq a_{w,v}^{(n,m)}}
       \widehat J_{v,a}X_v\widehat J_{v,a}^{T}\\
 &\quad-\sum_v\sum_{a_{w,v}^{(n,m)}<b\leq t_wc_{m,q}(v)}
       \widehat K_{v,b}X_v\widehat K_{v,b}^{T}.
 \end{aligned}
\]
For each $v$, the two sums contain altogether
$|a_{w,v}^{(n,m)}-t_wc_{m,q}(v)|$ terms.
Therefore,
\[
  \rk\bigl((\phi_{n,m}(X)-\sigma_{m,n,q}^{U}\rho_{m,q}(X))_w\bigr)\leq\sum_v|a_{w,v}^{(n,m)}-t_wc_{m,q}(v)|\rk(X_v).
\]
The flattened matrix $X_v$ has size $dp_m(v)\times dp_m(v)$, so
$\rk(X_v)\leq dp_m(v)$.  Therefore
\[
 \begin{aligned}
 \frac1{p_n(w)}\sum_v
 |a_{w,v}^{(n,m)}-t_wc_{m,q}(v)|\rk(X_v)
 &\leq d\sum_v
 \left|\frac{p_m(v)a_{w,v}^{(n,m)}}{p_n(w)}
       -\frac{p_m(v)t_wc_{m,q}(v)}{p_n(w)}\right|\\
 &=d\sum_v
       |\beta_{w,v}^{(n,m)}-\theta_w\lambda_{m,q}(v)|.
 \end{aligned}
\]
The last equality uses the definitions of $\beta$, $\theta_w$, and
$\lambda_{m,q}(v)$.

The triangle inequality gives
\[
 |\beta_{w,v}^{(n,m)}-\theta_w\lambda_{m,q}(v)|
 \leq|\beta_{w,v}^{(n,m)}-\alpha_m(v)|
       +|\alpha_m(v)-\lambda_{m,q}(v)|+(1-\theta_w)\lambda_{m,q}(v).
\]
The last term is nonnegative since $0\leq\theta_w\leq1$ and
$\lambda_{m,q}(v)\geq0$.

Summation over $v$ yields
\[
 \begin{aligned}
 \sum_v|\beta_{w,v}^{(n,m)}-\theta_w\lambda_{m,q}(v)|
 &\leq\|\beta_w^{(n,m)}-\alpha_m\|_1
       +\eta_{m,q}+(1-\theta_w)(1-\eta_{m,q})\\
 &\leq\|\beta_w^{(n,m)}-\alpha_m\|_1
       +\eta_{m,q}+(1-\theta_w).
 \end{aligned}
\]
The first inequality uses
$\sum_v|\alpha_m(v)-\lambda_{m,q}(v)|=\eta_{m,q}$ and
$\sum_v\lambda_{m,q}(v)=1-\eta_{m,q}$; the second uses
$0\leq\eta_{m,q}\leq1$.
Multiplying by $\alpha_n(w)$ and summing over $w\in V_n$ proves
\eqref{eq:parameter-first-square}.

Finally, apply \eqref{eq:return-map-composition} to
$Y=(y_{ij})\in\Mat_d(\Mat_q(R))$.
In the $(i,j)$-entry, the difference from $\jmath_{q,r}(Y)$ has its first $D_{n,q,r}$ copies equal to zero and
its remaining $r-D_{n,q,r}$ copies equal to $y_{ij}$.
The same coordinate permutation on rows and columns, ordering coordinates
first by the copy index and then by the matrix index, expresses the
difference as a block diagonal matrix with $r-D_{n,q,r}$ copies of $Y$ and a zero block.
Thus
\[
\rk_{qr}\bigl(
   \jmath_{q,r}(Y)-
   \rho_{m,n,q,r}^{U,P}\sigma_{m,n,q}^{U,P}(Y)
 \bigr)
 =\frac{r-D_{n,q,r}}r\rk_q(Y).
\]
Moreover,
\[
1-\frac{D_{n,q,r}}r
 =1-\sum_{w\in V_n}\lambda_{n,qr}(w)\theta_w
 =\eta_{n,qr}
   +\sum_{w\in V_n}\lambda_{n,qr}(w)(1-\theta_w)
 \leq\eta_{n,qr}+\tau_{n,q},
\]
because $\lambda_{n,qr}(w)\leq\alpha_n(w)$.
Since $\rk_q(Y)\leq d$, this proves
\eqref{eq:parameter-second-square}.
\end{proof}

\end{document}
